\section{A052392: a subdiagonal of Kimberling's triangle}
\label{sec:a052392}

\subsection{Source statement and indexing}

OEIS A052392 defines $a_n=T(2n+1,n)$, starting with $a_0=1$.
Robert Israel's November 26, 2018 contribution records an empirical
generating function and an empirical order-four recurrence. We prove both
as one source family. The algorithm attached to that entry defines
$T(r,0)=T(r,r)=1$, $T(r,k)=0$ for $k>r$, and $T(2,1)=3$; its remaining
entries satisfy
\begin{equation}\label{eq:a052392-triangle}
 T(r,k)=T(r-1,k-1)+2T(r-2,k-1)+T(r-1,k)
 \qquad(r\ge3,\ 0<k<r).
\end{equation}
The exceptional value $T(2,1)=3$ is essential. Applying the interior rule
there would instead produce $4$ and a different sequence.

The associated entry A054120 originally presents this triangle through
strip sums. Russell Jay Hendel's June 2, 2015 comment in that entry already
proves \eqref{eq:a052392-triangle} from the strip-sum definition and
explains its exceptional initial entry. That equivalence is credited
prior work. Our formal definition follows the literal algorithm attached
to A052392, rather than defining the sequence by the formula to be proved.

\begin{theorem}\label{thm:a052392-main}
Let $a_n=T(2n+1,n)$ with the preceding boundary conditions and recurrence.
In the formal power-series ring $\mathbb Q[[x]]$,
\begin{equation}\label{eq:a052392-gf}
 A(x):=\sum_{n\ge0}a_nx^n
 =\frac{2x^2-3x+1}{2x\sqrt{1-8x+4x^2}}+\frac{x-1}{2x},
\end{equation}
where the square root has constant coefficient $1$. The numerator of
the combined expression vanishes at zero, so the displayed expression
is a power series. Furthermore, for every $n\ge0$,
\begin{equation}\label{eq:a052392-rec}
\begin{split}
 8na_n-(32+28n)a_{n+1}+(78+30n)a_{n+2}\\
 {}-(45+11n)a_{n+3}+(5+n)a_{n+4}=0.
\end{split}
\end{equation}
\end{theorem}

\subsection{An independent array representation}

Define the weighted Delannoy array by
\[
 D(0,n)=D(m,0)=1,\qquad
 D(m+1,n+1)=D(m,n+1)+D(m+1,n)+2D(m,n).
\]
This definition determines every entry by induction on $m+n$ and gives
$D(m,n)=D(n,m)$. Write $S(m,n)=T(m+n,n)$. Thus $S$ has the same boundary
values and the same interior recurrence, except that $S(1,1)=3$.

\begin{lemma}\label{lem:a052392-difference}
For all $m,n\ge0$,
\[
 S(m+1,n+1)=D(m+1,n+1)-D(m,n).
\]
\end{lemma}
\begin{proof}
At $(m,n)=(0,0)$ the right side is $4-1=3$. Along the first interior
row, $D(1,n)=3n+1$, whereas $S(1,n)=3n$ for $n\ge1$; these identities
follow immediately by induction from the boundary rule. Symmetry gives
the first interior column. At all remaining positions, subtracting the
recurrences for $D(m+1,n+1)$ and $D(m,n)$ gives exactly the recurrence
for $S(m+1,n+1)$. Induction on $m+n$ now proves the result.
\end{proof}

Let
\[
 d_n=D(n,n),\qquad b_n=D(n+1,n).
\]
Then $d_0=b_0=a_0=1$, $d_1=4$, and Lemma
\ref{lem:a052392-difference} gives
\begin{equation}\label{eq:a052392-a-b}
 a_{n+1}=b_{n+1}-b_n\qquad(n\ge0).
\end{equation}
By symmetry and the array recurrence,
\begin{equation}\label{eq:a052392-d-b}
 d_{n+1}=2b_n+2d_n\qquad(n\ge0).
\end{equation}

\subsection{The diagonal recurrence}

\begin{lemma}\label{lem:a052392-balance}
For every $m,n\ge0$,
\begin{equation}\label{eq:a052392-balance}
 (m+1)D(m+1,n)-(n+1)D(m,n+1)=2(n-m)D(m,n).
\end{equation}
\end{lemma}
\begin{proof}
Let $E(m,n)$ be the left side minus the right side of
\eqref{eq:a052392-balance}. On the axes, the formula
$D(1,n)=3n+1$ and symmetry give $E(0,n)=E(m,0)=0$.
Expanding three instances of the defining recurrence for $D$ gives
\[
 E(m+1,n+1)=E(m,n+1)+E(m+1,n)+2E(m,n).
\]
For an explicit check, subtract the right side from the left; the
remaining expression is
\[
 2(m-n)\bigl(D(m+1,n+1)-D(m,n+1)-D(m+1,n)-2D(m,n)\bigr)=0.
\]
Induction on $m+n$ therefore gives $E(m,n)=0$ everywhere.
\end{proof}

Apply \eqref{eq:a052392-balance} at $(m,n)=(n+1,n)$, and use
\[
 b_{n+1}=D(n+1,n+1)+D(n+2,n)+2D(n+1,n).
\]
The result is
\begin{equation}\label{eq:a052392-brec}
 (n+2)b_{n+1}=(2n+3)d_{n+1}+2(n+1)b_n.
\end{equation}
Eliminating $b_n,b_{n+1}$ with \eqref{eq:a052392-d-b} yields
\begin{equation}\label{eq:a052392-drec}
 (n+2)d_{n+2}=4(2n+3)d_{n+1}-4(n+1)d_n\qquad(n\ge0).
\end{equation}
Likewise, \eqref{eq:a052392-a-b} and \eqref{eq:a052392-d-b} give
\begin{equation}\label{eq:a052392-source-central}
 2a_{n+1}=d_{n+2}-3d_{n+1}+2d_n\qquad(n\ge0).
\end{equation}
Every relation so far is an identity for the independently defined arrays.

\subsection{The generating function}

Put $D(x)=\sum_{n\ge0}d_nx^n$ and $Q(x)=1-8x+4x^2$.
Equation \eqref{eq:a052392-drec}, together with $d_0=1$ and $d_1=4$,
is equivalent coefficient by coefficient to
\begin{equation}\label{eq:a052392-ode}
 QD'=(4-4x)D.
\end{equation}
Indeed, the coefficient of $x^j$ is
$(j+1)d_{j+1}=(8j+4)d_j-4jd_{j-1}$ for $j\ge1$,
and the constant coefficient says $d_1=4d_0$.
Formal differentiation now gives
\[
 (QD^2)'=(-8+8x)D^2+2D\,QD'=0.
\]
Over $\mathbb Q[[x]]$, a series with zero derivative is constant, and the
constant coefficient here is $1$. Thus
\begin{equation}\label{eq:a052392-square}
 QD^2=1,\qquad D(0)=1.
\end{equation}
These conditions characterize $D$ uniquely: if $QE^2=1$ and $E(0)=1$,
then $Q\ne0$ gives $(E-D)(E+D)=0$. The ring is a domain and $E+D$ has
constant coefficient $2$, so $E=D$. Consequently $D=1/\sqrt Q$ with
the normalized branch. This argument proves existence and uniqueness;
it does not assume a square-root coefficient expansion.

Multiply \eqref{eq:a052392-source-central} by $x^{n+2}$ and sum over
$n\ge0$, using $a_0=d_0=1$ and $d_1=4$. One obtains
\begin{equation}\label{eq:a052392-cleared}
 2xA(x)=(1-3x+2x^2)D(x)-1+x
       =(1-x)\bigl((1-2x)D(x)-1\bigr).
\end{equation}
Substitution of $D=1/\sqrt Q$ gives exactly
\eqref{eq:a052392-gf}. In particular the right side of
\eqref{eq:a052392-cleared} has zero constant coefficient. Multiplication
by $2x$ is injective in $\mathbb Q[[x]]$, so the cleared identity loses
no information about $A$.

\subsection{The exact source recurrence}

For clarity, an explicit algebraic certificate gives the recurrence
without relying on a guessed differential equation for $A$. Define
\[
 R_j=(j+2)d_{j+2}-4(2j+3)d_{j+1}+4(j+1)d_j.
\]
By \eqref{eq:a052392-drec}, every $R_j=0$.
Let $L_n$ denote the left side of \eqref{eq:a052392-rec}.
Substitute \eqref{eq:a052392-source-central} into the five terms of
$2L_{j+1}$. Collecting coefficients of $d_j,\ldots,d_{j+6}$ gives the
polynomial identity
\[
 2L_{j+1}=4R_j-12R_{j+1}+13R_{j+2}-6R_{j+3}+R_{j+4}=0
 \qquad(j\ge0).
\]
Hence $L_n=0$ for every $n\ge1$. Finally the source array gives
$(a_0,a_1,a_2,a_3,a_4)=(1,6,39,261,1779)$, and
\[
 L_0=-32\cdot6+78\cdot39-45\cdot261+5\cdot1779=0.
\]
This proves \eqref{eq:a052392-rec} at every source index and completes
the theorem.

\subsection{Formal coverage and provenance}

The Lean development defines the source array and the weighted Delannoy
array independently. It proves their relation, the balance identity,
the diagonal recurrence, the normalized formal reciprocal-square-root
identity and uniqueness, the cleared generating function, and the exact
order-four recurrence for every natural index. The single audited
endpoint is \texttt{A052392.full\_resolution}. Its dependencies are only
the standard Lean axioms \texttt{propext}, \texttt{Classical.choice}, and
\texttt{Quot.sound}. All four local modules are explicit default build
roots; the compiler and mathlib revision are pinned.

The frozen sources are OEIS A052392 revision 12 (April 24, 2025) and
A054120 revision 27 (January 26, 2022), retrieved from immutable
OEIS-data commits. Weighted Delannoy methods and generalized Pascal
triangles are classical; see Paul Barry, \emph{On Integer-Sequence-Based
Constructions of Generalized Pascal Triangles}, Journal of Integer
Sequences \textbf{9} (2006), Article 06.2.4. The exact and semantic
prior-work check is documented separately, including known source
implementations and the current ProveIt and trureturing repositories.
No proof of the precise A052392 empirical assertions was found within
that bounded search. We assert resolution of this recorded source
family, without asserting worldwide priority for its underlying methods.
